% Compiler avec XeLaTeX ou LuaLaTeX
\documentclass[11pt,a4paper]{article}

\usepackage{amsmath,amssymb}
\usepackage{fontspec}
\usepackage{unicode-math}
\usepackage{polyglossia}
\setmainlanguage{french}
\usepackage{xcolor}
\usepackage{fancyhdr}
\usepackage[margin=2cm,headheight=16pt,headsep=22pt]{geometry}
\usepackage{enumitem}

% Police principale et police mathématique : Libertinus
\IfFontExistsTF{Libertinus Serif}
  {\setmainfont{Libertinus Serif}}
  {\IfFontExistsTF{Latin Modern Roman}{\setmainfont{Latin Modern Roman}}{\setmainfont{lmroman10-regular.otf}}}

\IfFontExistsTF{Libertinus Math}
  {\setmathfont{Libertinus Math}}
  {\IfFontExistsTF{Latin Modern Math}{\setmathfont{Latin Modern Math}}{\setmathfont{latinmodern-math.otf}}}

\linespread{1.07}
\setlength{\parindent}{0pt}
\AtBeginDocument{\setlength{\parindent}{0pt}}
\setlength{\parskip}{0.35em}

% En-tête
\pagestyle{fancy}
\fancyhf{}
\lhead{Nom de l'enseignant}
\rhead{classe à préciser}
% Pagination : « Page X / Y » centrée en pied de page
\usepackage{lastpage}
\cfoot{Page \thepage\ / \pageref{LastPage}}
\renewcommand{\headrulewidth}{0.5pt}
% Trait de séparation éloigné du texte de l'en-tête (modifier 5pt pour ajuster)
\renewcommand{\headrule}{\vskip 5pt\hrule width\headwidth height\headrulewidth \vskip-\headrulewidth}

% Configuration des listes
\setlist[enumerate,1]{itemsep=0.45em, topsep=0.3em, label=\textbf{\arabic*.}}
\setlist[enumerate,2]{itemsep=0.25em, topsep=0.2em, label=\textbf{\alph*.}}

% Commande exercice
\newcommand{\exo}[2]{%
  \vspace{1.5em}%
  \noindent{\large\textbf{Exercice #1. -- #2}}%
  \par\vspace{0.35em}%
  \hrule height 0.5pt%
  \vspace{0.7em}%
}

\newcommand{\partie}[1]{%
  \vspace{0.9em}%
  \noindent\textbf{#1}%
  \par\vspace{0.3em}%
}

\begin{document}

\begin{center}
  {\LARGE \textbf{Titre à préciser}}\\[0.5em]
  {\large Classe à préciser}
\end{center}

% ---------------------- EXERCICE 1 ----------------------
\exo{1}{Titre de l'exercice}

\partie{Partie A -- Titre de la partie}

Énoncé introductif.
\begin{enumerate}
  \item Question.
  \item Question.
  \begin{enumerate}
    \item Sous-question.
    \item Sous-question.
  \end{enumerate}
\end{enumerate}

\partie{Partie B -- Titre de la partie}

\begin{enumerate}
  \item Question.
\end{enumerate}

% ---------------------- EXERCICE 2 ----------------------
\exo{2}{Titre de l'exercice}

Énoncé introductif.
\begin{enumerate}
  \item Question.
  \item Question.
  \begin{enumerate}
    \item Sous-question.
    \item Sous-question.
  \end{enumerate}
\end{enumerate}

\end{document}

